% !TEX root = ../nc-cubic.tex
\subsubsection{境界の低次積・二次関係の飽和と規格化周期}
\label{subsubsec:nfm-boundary-products-period}

本節では、元の trait 上で指定周期を保った低次積を計算する。
全次元の orbit Hom は平坦である一方、特殊 fibre の二次積は階数を三つ失う。
この三つの関係を飽和すると平坦な表示が回復し、中心三次の最初の項と
滑らかな potential 三次の同型延長から、完備局所周期位相を決定できる。
固定 MF・window・一般点生成性・飽和性の入力は引き続き前提とする。

\begin{lemma}[全次元の右窓と点線束の section module]
\label{lem:nfm-boundary-right-window-module-alln}
$R$ を本境界の DVR、$S=R[X,Y,Z,u_1,\ldots,u_r]$、
$A=S/(F_R)$、$r=n-1\ge2$ とし、既存の右射影で
$Z_R=j_R^!i_*\mathcal M[1]$ を定める。
$N_0=\operatorname{coker}(L:R[X,Y,Z](-1)^3\to R[X,Y,Z]^3)$ を
中心変数が零として作用する $A$-加群と見なす。
このとき window $i=1$ の enhanced functor により
\[
 \Psi_1Z_R\simeq N_0(-1)[1]\quad\text{in }D_{\mathrm{sg}}^{\mathrm{gr}}(A),
 \qquad\Psi_1=\pi_{\mathrm{sg}}R\Gamma_{\ge1}
\]
である。これは閉 fibre と一般 fibre の両方で、
内部 shift と方程式 $F_R$ の指定 MF 周期を保つ同定である。
\end{lemma}

\begin{proof}
適用する窓定理の環論的条件を確認する。
$S$ は regular DVR 上の多項式環、$F_R$ は正則な斉次三次である。
従って $A$ は noetherian かつ Gorenstein であり、graded duality により
$R\operatorname{Hom}_A(R,A)\simeq R(r)[1-(n+2)]$ となる。
$A_0=R$ は有限 global dimension を持ち、可換 noetherian graded 環なので
条件 $\chi$ を満たす。
Burke--Stevenson\cite[Definition~5.1, Proposition~4.18, Theorem~6.4(1)]{NFMBurkeStevenson}
を parameter $a=r$、window $i=1$ で適用すると
\[
 D^b(\operatorname{coh}\mathcal Y)
 =\langle\mathcal O(-r),\ldots,\mathcal O(-1),\mathcal B_1\rangle
\]
を得る。相対双対性で
$\mathcal B_1\cap\operatorname{Perf}(\mathcal Y)=\mathcal A_R$ である。
grade truncation、導来双対、自由 summand の除去という証明中の構成を
dg quotient で行い、$\Psi_1$ と窓の enhanced 同定を固定する。
ここでは単なる triangulated 同値から高次積の整合を推論しない。

$1\le j\le r$ と $t\ge1$ に対し $t-j\ge1-r$ である。
三次超曲面の完全列から中間 cohomology は零であり、最高次も
$H^n(\mathcal O(t-j))\simeq H^0(\mathcal O(j-t-r))^*=0$ である。
従って
\[
 R\Gamma_{\ge1}\mathcal O(-j)=A(-j)
\]
であり、この対象は特異圏で零となる。
右射影の有限個の mutation triangles を $\Psi_1$ で移すと
$\Psi_1(j_Rj_R^!i_*\mathcal M)=\Psi_1(i_*\mathcal M)$ となる。

$t\ge1$ では $H^1(\mathcal C,\mathcal M(t))=0$ であり、
線型表示の完全列から section module は
\[
 R\Gamma_{\ge1}(i_*\mathcal M)
 =\bigoplus_{t\ge1}H^0(\mathcal C,\mathcal M(t))=N_0(-1)
\]
となる。よって所要の式が従う。
これらの section modules と各有限 cone は $R$ 上で定義され、
必要な cohomology は自由なので閉／一般底変更と両立する。
これは任意の特異圏についての包括的な底変更定理ではない。

隣接する窓の grade truncation と自由 summand の除去を比較すると、
内部 shift $(-1)$ は右 mutation $R_{\mathcal O(-1)}(-1)$ に対応する。
MF 側では同じ曲率 $F_R$ に対する $(-3)=[-2]$ を用いる。
この実際の周期を窓へ移しているので、scalar の規格化変更はない。
共通の内部 shift $(-1)$ と cohomological shift $[1]$ は、
標準の符号付き dg 同型によって endomorphism 代数を変えない。
特に単位を保ち、周期の cohomological degree は二なので周期の符号も変わらない。
\end{proof}

\begin{lemma}[零曲率の中心 Koszul 因子の dg 分離]
\label{lem:nfm-boundary-koszul-dg-separation}
$T=k[u_1,\ldots,u_r]$、
$K_T=(T\otimes\Lambda(\epsilon_1,\ldots,\epsilon_r),d=\sum u_i\iota_i)$ とし、
$\deg\epsilon_i=(-1,1)$ とする。
ここで二次数は cohomological degree と Adams degree である。
縮約 $\iota_i$ を $e_i$ の像とする写像
\[
 \Lambda_k(e_1,\ldots,e_r)\longrightarrow\operatorname{End}_T(K_T),
 \qquad\deg e_i=(1,-1)
\]
は unital bigraded dg 代数の quasi-isomorphism である。
active 三次 $H$ の有限 graded MF $P$ に対し、tensor は
\[
 \operatorname{End}(P)\otimes\Lambda_k(e_1,\ldots,e_r)
 \longrightarrow\operatorname{End}(P\widehat\otimes K_T)
\]
という dg 代数の quasi-isomorphism を与え、実際の曲率 $H$ の MF 周期を厳密に保つ。
\end{lemma}

\begin{proof}
縮約は互いに反交換し、$d$ と超可換するので、表示した写像は dg 代数の射である。
augmentation $K_T\to k$ への postcomposition は
$\operatorname{End}_T(K_T)\to\operatorname{Hom}_T(K_T,k)$ という複体の射を与える。
$K_T$ は有限自由分解なので、この射は quasi-isomorphism である。
右辺の微分は零で、縮約の像はその外積基底を与える。
従って元の dg 代数の射も quasi-isomorphism である。
finite MF と有限 Koszul 複体の tensor の Hom は、符号付き tensor Hom に一致する。
$k$ は体なので tensor は quasi-isomorphism を保つ。MF 周期は $P$ の周期と
$K_T$ の恒等写像の tensor であり、scalar 変更はない。
\end{proof}

\begin{remark}[高次積を分離する際の意味]
active な厳密単位的 minimal $A_\infty$ 模型を、次数付き可換 dg 代数
$\Lambda(e_1,\ldots,e_r)$ で係数拡大すれば、
\[
 m_k(a_1\otimes\lambda_1,\ldots,a_k\otimes\lambda_k)
 =\pm m_k^{\mathrm{act}}(a_1,\ldots,a_k)
            \otimes\lambda_1\cdots\lambda_k
\]
という full 模型を取れる。
とくに $k\ge3$ で入力に純粋な中心生成元 $1\otimes e_i$ があれば積は零である。
任意の別の contraction で同じ係数公式が文字通り成立するという主張ではなく、
指定周期を保つ $A_\infty$ 同値の下でこの模型を選べるという意味である。
\end{remark}

\begin{proposition}[分離境界の active 二次・三次積は消える]
\label{prop:nfm-separated-boundary-active-products}
$L=L_0$、$Q_a=-b\operatorname{adj}L$ とし、
\[
 \mathsf D=\begin{pmatrix}0&L\\Q_a&0\end{pmatrix},\qquad
 \mathsf D^2=bXYZ I_6
\]
を active MF の微分とする。
その Adams graded endomorphism dg 代数を $\mathscr E_a$ と書く。
$E_{ij}$ を三行三列の行列単位とすると、
$H^{1,-1}(\mathscr E_a)$ は次の三つの cocycle の class を基底に持つ：
\[
 \begin{aligned}
 \phi_a&=\begin{pmatrix}0&-E_{21}\\bYE_{31}&0\end{pmatrix},\quad
 \phi_b=\begin{pmatrix}0&-E_{31}\\b(XE_{21}+2ZE_{31})&0\end{pmatrix},\\
 \phi_c&=\begin{pmatrix}0&-E_{22}\\b(XE_{31}+ZE_{33})&0\end{pmatrix}.
 \end{aligned}
\]
さらに
\[
 H^{2,-2}(\mathscr E_a)=k\{X\operatorname{id},Y\operatorname{id},Z\operatorname{id}\},
 \qquad H^{1,-2}(\mathscr E_a)=0,
 \qquad H^{2,-3}(\mathscr E_a)=k\operatorname{id}.
\]
ここで最後の identity は指定 MF 周期の class である。
この三次周期で正規化した minimal 模型において
\[
 m_2\big|_{(H^{1,-1})^{\otimes2}}=0,\qquad
 m_3\big|_{(H^{1,-1})^{\otimes3}}=0.
\]
後者の消滅も minimal 模型の gauge に依存しない。
\end{proposition}
\begin{proof}
微分は $\delta(f)=\mathsf Df-(-1)^{|f|}f\mathsf D$ である。
$(1,-1)$ 次数の cochain は、右上に定数行列 $C$、左下に一次行列を持ち、
計三十六係数を持つ。左下を $bD$ と書けば cocycle 条件は
\[
 LD-C\operatorname{adj}L=0,\qquad
 DL-\operatorname{adj}L\,C=0.
\]
この百八行三十六列の係数行列は階数三十三であり、
三十三次小行列の行列式は $-256p^2$ である。
表示した三つの独立な解があるので階数の上界も一致する。
この次数の boundaries は次数の理由で零である。

$(2,-2)$ 次数の cochain は二つの一次対角 block $(A,B)$ であり、
条件は $LB-AL=0$、$\operatorname{adj}L\,A-B\operatorname{adj}L=0$ である。
百四十四行五十四列の係数行列に四十二次 minor $-64p^3$ がある。
九次元の boundaries は定数行列 $C$ による $(LC,CL)$ で、
その独立性は九次 minor $-8$、または $L$ の一般点での可逆性から従う。
さらに $X\operatorname{id},Y\operatorname{id},Z\operatorname{id}$ が cocycle である。
これらの線型結合 $\ell\operatorname{id}$ が boundary なら、
$LC=\ell I_3$ から $\ell^3=-XYZ\det C$ となる。
三つの相異なる一次素因子を比較すると $\ell=0$、従って $C=0$ である。
これで kernel は十二次元、cohomology は表示した三次元となる。

$(1,-2)$ の cochain は左下の定数行列だけであり、その微分は $(LC,CL)$
なので cocycle は零である。
$(2,-3)$ は二つの定数対角 block である。
対応する八十一行十八列の係数行列は十七次 minor $-8p^2$ を持ち、
identity が kernel を張る。この次数の boundaries は零である。
これらの minor の添字と係数行列の生成式は、自己完結した検証器
\texttt{computations/verify\_boundary\_low\_products.py} に含めた。

次に積を計算する。九つの順序付き積のうち、唯一零でないものは
\[
 \phi_c\phi_b=\delta T,\qquad
 T=\begin{pmatrix}0&0\\-bE_{31}&0\end{pmatrix},
 \qquad \deg T=(1,-2)
\]
である。その他の八積は cochain として零である。
従って $m_2$ は零となる。
しかも全 $i\in\{a,b,c\}$ に対し
\[
 \phi_iT=T\phi_i=0
\]
が cochain として成り立つ。
$H^{1,-2}=0$ なのでこの boundary の homotopy は一意であり、
$h(\delta T)=T$ とする contraction を選べる。
この contraction を他の次数と strict unit に両立して延ばす。
arity 三の transfer 式の二つの項は上の二つの積なので、共に零である。
これで指定周期成分を含む $m_3$ の消滅を得る。

一次成分が恒等の gauge による変化では、
$f_2(H^{1,-1},H^{1,-1})\subset H^{1,-2}=0$ であり、
残る項も $m_2$ の上の制限が零なので消える。
従って三次消滅は gauge に依存しない。
\end{proof}

\begin{proposition}[split 境界模型の全 orbit Hom 平坦性と二次積の階数落ち]
\label{prop:nfm-boundary-all-orbit-flat-rank-drop}
$n\ge3$、$r=n-1$、$N=n+2$ とする。
pole 消去を厳密な不等式で行い、特殊方程式が $F_0=bXYZ$ となる模型を取る。
右窓の対象 $Z_R$ と、その実際の MF 周期を固定する。
一般 fibre については既存の固定 Tate--Koszul 入力を仮定する。
このとき全ての整数 $i,q$ に対し、
\[
 S_R^{i,q}=H^i\operatorname{RHom}(Z_R,\rho^qZ_R)
\]
は $R$-自由で、閉 fibre と一般 fibre への底変更に両立し、
その階数は $\binom N{3i+2q}$ である。範囲外の二項係数は零とする。
しかし二次積
\[
 \mu_R:(S_R^{1,-1})^{\otimes2}\longrightarrow S_R^{2,-2}
\]
の特殊 fibre の階数は $3r+\binom r2=\binom N2-3$ である。
その一般 fibre は全射なので、$\operatorname{coker}\mu_R$ は
正整数 $a_1,a_2,a_3$ に対する
$\bigoplus_{i=1}^3R/(z^{a_i})$ という torsion 加群である。
指数は後の命題\ref{prop:nfm-boundary-first-jet-products}で全て一と決定する。
\end{proposition}

\begin{proof}
補題\ref{lem:nfm-boundary-right-window-module-alln}により、$Z_R$ は点線束の section module $N_0(-1)[1]$ の
stabilization で表され、その同定は指定周期を保つ。
特殊 fibre の MF は active MF と $K_T$ の tensor なので、補題\ref{lem:nfm-boundary-koszul-dg-separation}が適用される。

active 部では、$q=0$ の cohomology は次数零と一に一つずつ、
$q=-1$ は次数一に三つ、$q=-2$ は次数二に三つである。
最初は $H^*(\mathcal O_C)$、二番目は $H^*(\mathcal O_C(-1))$、
三番目は右 mutation の triangle と、余因子 pairing による全射
$\mu:H^0(M^{-1}(1))\otimes H^0(M(1))\to H^0(\mathcal O_C(2))$
の双対の完全列から従う。これは $\mu^*:k^6\hookrightarrow k^9$ の cokernel が三次元という計算であり、他の cohomology は零である。
$\sigma_C^3=[-2]$ により、全 active orbit cohomology の二次数ベクトル空間は
$\Lambda(k^3)\otimes k[t,t^{-1}]$、$\deg t=(2,-3)$ と一致する。
これは積の同定ではない。
中心外積因子を tensor すると全特殊 fibre のベクトル空間は
$\Lambda(k^N)\otimes k[t,t^{-1}]$ となる。
一般 fibre は固定入力の $B^!\otimes K[t,t^{-1}]$ であり、
$H_{B^!}(s)=(1+s)^N$ なので同じ二次数次元を持つ。
各固定 $q$ の Hom は proper perfect な有限複体であるから、
DVR の普遍係数完全列により、その cohomology は自由となる。
$\deg t=(2,-3)$ から外積次数は $3i+2q$ であり、表示した階数式を得る。

特殊 fibre の次数一は active 三次元と中心 $r$ 次元に分かれ、
次数二は active 三次元、混合 $3r$ 次元、中心外積 $\binom r2$ 次元に分かれる。
active 二次積は命題\ref{prop:nfm-separated-boundary-active-products}により零で、残る二つの summand へは
tensor 積と外積が全射である。従って階数は $3r+\binom r2$ である。
一般 fibre の二次積は quadratic Koszul algebra の双対の次数二への積なので全射である。
従って cokernel は torsion であり、その剰余次元は三である。
DVR 上の構造定理から、非零な巡回 torsion 因子の個数は正確に三と分かる。
\end{proof}

\begin{corollary}[全次元の integral な低次積と原始的 target 周期]
\label{cor:nfm-boundary-alln-integral-products}
前命題の模型では、全 $n\ge3$ について、実際の MF 周期を固定した
strictly unital、Adams graded な $R$ 上の minimal $A_\infty$ 模型を取れる。
$N=n+2$ とすると
\[
 (S_R^{1,-1})^*\simeq R^N,\qquad
 (S_R^{2,-2})^*\simeq R^{\binom N2},\qquad
 S_R^{2,-3}=Rt_p
\]
であり、$t_p$ は原始的である。
\end{corollary}
\begin{proof}
全 cohomology の自由性が得られたので、
系\ref{cor:nfm-boundary-threefold-integral-low-products}の単一の enhanced MF 模型上の
contraction と finite tree の構成を全 $n$ に適用できる。
単位と周期を固定し、三つの Adams 周期剰余で分裂を選べばよい。
周期は $H^0\operatorname{End}(Z_R)=R\operatorname{id}$ から来るので原始的である。
ここでは target の周期と回復中心三次の原始性を同一視しない。
\end{proof}

\begin{lemma}[full 模型と active 計算の一次 jet の一致]
\label{lem:nfm-boundary-full-active-first-jet}
中心座標の pole 消去をさらに強めて
\[
 F_R\equiv J_R\pmod{z^2},\qquad F_R-J_R=\sum u_iG_i,\quad G_i\in z^2R[X,Y,Z,u]_2
\]
とできる。このとき、上の $R$ 上の minimal 模型を、$R/(z^2)$ 上で
active な相対行列 $L(q=2+z)$ の模型と $\Lambda(e_1,\ldots,e_r)$ との
係数拡大に一致するように取れる。この同定は実際の指定周期を保つ。
\end{lemma}
\begin{proof}
命題\ref{prop:nfm-hesse-boundary-integral-residual-model}の指数を、全係数について
$v_z(c)+Nj\ge2$ を満たすまで増やせばよい。一般 fibre の位相は変わらない。
補題\ref{lem:nfm-boundary-line-mf}の actual MF は、$R/(z^2)$ 上で active MF と
零曲率の中心 Koszul 複体の tensor に厳密に一致する。
補題\ref{lem:nfm-boundary-koszul-dg-separation}の外積代数からの写像は、
この非被約基底でも quasi-isomorphism である。実際、変数 $u_i$ は任意の基底上で
正則列であり、有限自由 Koszul 分解の augmentation に Hom を施す同じ証明が使える。
各内部次数でその cone は有界な自由複体で分裂するため、active Hom との tensor にも
quasi-isomorphism が保たれる。

全 cohomology が $R$-自由なので、各 Hom 複体の cycles と boundaries は自由で、
それらを定める短完全列は分裂する。この分裂の $R/(z^2)$ 上での任意の指定を、
自由基底を持ち上げて $R$ 上の分裂へ延ばせる。
この操作を同じ MF の三つの周期剰余で行い、単位を固定して周期的に反復する。
まず tensor dg 包含の像への retract を、その像上で homotopy が零となるように取り、
次に active contraction を外積代数で係数拡大する。
active 側の厳密単位性により、純粋な中心入力を含む三次積も零となる。
従って係数拡大による contraction と一致する一次 jet を持つ contraction を選べる。
有限 tree 式もこの底変更と可換なので、二次・三次積の一次 jet が一致する。
これは任意に選んだ別々の Hom 模型の積を貼り合わせる議論ではない。
\end{proof}

\begin{proposition}[消えた三つの二次関係は一次で回復する]
\label{prop:nfm-boundary-first-jet-products}
境界パラメータを $\varepsilon=q-2$ と書き、$d=p^3+8$、$b_0=d/3$ と置く。
補題\ref{lem:nfm-hesse-boundary-explicit-line}の実際の相対行列
$L(\varepsilon)$ と $-b_0\operatorname{adj}L(\varepsilon)$ を用いる。
命題\ref{prop:nfm-separated-boundary-active-products}の三つの active cocycle を $\phi_a,\phi_b,\phi_c$ と書き、
その定数上右 block の $(2,1),(3,1),(2,2)$ 成分を固定して一次まで持ち上げる。
$H^{2,-2}$ の基底を $(X\operatorname{id},Y\operatorname{id},Z\operatorname{id})$
とすると、列を $(aa,ab,ac,ba,bb,bc,ca,cb,cc)$ の順に並べた積は
\[
 m_2=\frac{\varepsilon}{3}
 \begin{pmatrix}
 0&0&p^{-1}&0&p&0&0&0&0\\
 p^{-1}&0&0&0&0&0&0&1&0\\
 0&0&0&2p^{-1}&0&0&0&0&p/4
 \end{pmatrix}+O(\varepsilon^2).
\]
従って三つの active elementary divisor の付値は全て一である。
飽和した二次 relation の特殊 fibre は
\[
 ac+p^2b^2,\qquad a^2+p\,cb,\qquad ba+\frac{p^2}{8}c^2
\]
である。ここで $a,b,c$ は cohomology の双対生成元である。
生成元を $x=a$、$y=pb$、$z=pc/2$ と改めると、これらは
\[
 yx=-\frac p2z^2,\qquad zy=-\frac p2x^2,\qquad xz=-\frac p2y^2
\]
という元の active 境界代数の関係に一致する。
\end{proposition}
\begin{proof}
$L=L_0+\varepsilon L_1+O(\varepsilon^2)$ とし、
$\phi_i=\begin{psmallmatrix}0&C_i\\b_0D_i&0\end{psmallmatrix}$ と書く。
$C_i=C_{i,0}+\varepsilon C_{i,1}+O(\varepsilon^2)$、
$D_i=D_{i,0}+\varepsilon D_{i,1}+O(\varepsilon^2)$ とする。
$J=\operatorname{adj}L$ と置けば、閉性は
\[
 LD_i-C_iJ=0,\qquad D_iL-JC_i=0
\]
であり、その一次式は特殊 fibre と同じ階数三十三の線型系である。
指定した三成分の正規化を加えると解が一意となる。
相対 $L$ の有理式を微分して解くと
\[
 \begin{aligned}
 C_{a,1}&=\frac1d\begin{pmatrix}-1&0&p^{-1}\\0&0&-2p^{-1}\\0&-2p^{-1}&0\end{pmatrix},
 \quad
 C_{b,1}=\frac1d\begin{pmatrix}p&p&0\\0&0&-2\\0&0&-2p^{-1}\end{pmatrix},\\
 C_{c,1}&=\frac1d\begin{pmatrix}-1&-1&0\\0&0&0\\0&-2p^{-1}&p/4\end{pmatrix}
 \end{aligned}
\]
を得る。また、積の計算に必要な下左 block の成分は
\[
 (D_{a,1})_{13}=-\frac{2Z}{pd},\qquad
 (D_{b,1})_{13}=-\frac{pX}{d},\qquad
 (D_{c,1})_{13}=-\frac{Y+2Z/p}{d}
\]
である。これらの式を含む全一次閉性は、有理係数の有限行列として検算できる。

次数 $(2,-2)$ の cocycle を $(A,B)$ とすると、その cohomology 座標は
\[
 ([X]A_{11},\ [Y]A_{11},\ [Z]A_{33})
\]
で与えられる。実際、boundary の上 block は $L_0H$ ($H$ 定数) なので、
$(L_0H)_{11}$ は $Z$ の倍数、$(L_0H)_{33}$ は $X,Y$ の線型結合であり、
この三成分は全 boundary を消す。また三つの scalar identity 上で標準基底となる。

唯一の非零な零次積 $\phi_c\phi_b=\delta T$ については、
$T$ の零次値を保ち、一次積から $\delta_1T$ を引いてからこの射影を用いる。
その他の積にはこの補正は不要である。
表示した $C_{i,1}$ と三つの $(D_{i,1})_{13}$ を代入すると上の行列となる。
三列 $(aa,ac,ba)$ の一次係数の determinant は
\[
 -\frac2{27p^3}\in k^\times
\]
である。従って、全係数が $\varepsilon$ で割れ、割った写像が全射なので、
Smith 付値はちょうど $(1,1,1)$ である。
双対の行を非零 scalar で規格化して、表示した三関係を得る。
最後の生成元変更は直接代入で確認できる。
\end{proof}

\begin{proposition}[飽和した二次関係が定める全次数の平坦族]
\label{prop:nfm-boundary-saturated-quadratic-flat}
上の積分模型で、中心を含む方程式係数が $z^2$ で割れるように
pole 消去の指数を取る。
$X_R=(S_R^{1,-1})^*$、$\mathcal R_2=\operatorname{im}(\mu_R^*)$、
\[
 \mathcal R_2^{\mathrm{sat}}
   =(\mathcal R_2\otimes_RK)\cap X_R^{\otimes2},
 \qquad
 B_R^{\mathrm{sat}}=T_R(X_R)/(\mathcal R_2^{\mathrm{sat}})
\]
と置く。
命題\ref{prop:nfm-boundary-first-jet-products}と full 一次 jet の一致により、三つの非単元 Smith 因子が全て $z$ であり、
閉 fibre の active 関係が
\[
 ac+p^2b^2=0,\qquad a^2+pcb=0,\qquad ba+(p^2/8)c^2=0
\]
となる。残りの関係は、補題\ref{lem:nfm-boundary-full-active-first-jet}により中心変数との可換関係である。
このとき $B_R^{\mathrm{sat}}$ の次数 $d$ の部分は有限自由で、階数は
$\binom{n+d+1}{d}$ である。
従って全次数で $R$-flat な quadratic algebra を得る。

また $B_R^{\mathrm{raw}}=T_R(X_R)/(\mathcal R_2)$ と置けば、
$B_R^{\mathrm{sat}}$ は $B_R^{\mathrm{raw}}$ をその $R$-torsion 全体で割ったものに等しい。
その torsion ideal は次数二の三つの元で生成され、$z$ で消える。
\end{proposition}

\begin{proof}
$\mathcal R_2^{\mathrm{sat}}$ は有限自由加群 $X_R^{\otimes2}$ 内の飽和部分加群なので、
その商は自由であり、部分加群の階数は一般 fibre と同じ $\binom{n+2}{2}$ である。
$X=a$、$Y=pb$、$Z=pc/2$ と置くと、active 関係は
\[
 ZY=-(p/2)X^2,\qquad XZ=-(p/2)Y^2,\qquad YX=-(p/2)Z^2
\]
となる。
$\upsilon=-p/2\in k^\times$ と置き、量子多項式環
\[
 T_{\upsilon}=k\langle X,Y,Z\rangle/
       (ZX-\upsilon XZ,\ XY-\upsilon YX,\ YZ-\upsilon ZY)
\]
と、その巡回自己同型 $\tau(X)=Z$、$\tau(Y)=X$、$\tau(Z)=Y$ を考える。
積 $a*b=a\tau^{\deg a}(b)$ による Zhang twist の二次関係は、ちょうど上の三式である。
これは tensor algebra の各次数で
$a_1\otimes\cdots\otimes a_d\mapsto
 a_1\tau(a_2)\cdots\tau^{d-1}(a_d)$ を用いても確認できる。
従って active 特殊代数の Hilbert 級数は $(1-t)^{-3}$ である。
中心変数 $r=n-1$ 個の多項式環を tensor するので、
全特殊代数の Hilbert 級数は $(1-t)^{-(n+2)}$ となる。

一方、一般 fibre の二次関係は固定 Tate--Koszul 入力による元の
quadratic AS-regular algebra $B_K$ の関係そのものである。
従って一般 fibre も Hilbert 級数 $(1-t)^{-(n+2)}$ を持つ。
各 $d$ で $(B_R^{\mathrm{sat}})_d$ は有限生成 $R$-加群である。
表示による代数の商は底変更に両立するので、その閉 fibre の次元も一般 rank も
$\binom{n+d+1}{d}$ である。
DVR 上の有限生成加群では、剰余次元は rank と非零巡回 torsion 因子数の和である。
従ってこの加群は torsion を持たず、有限自由である。
これが全次数について成立する。

$B_R^{\mathrm{raw}}\twoheadrightarrow B_R^{\mathrm{sat}}$ の核は
一般 fibre で零となるため $R$-torsion である。
逆に任意の torsion 元は torsion-free な標的へ零に写るので、核は torsion 全体である。
核は定義から $\mathcal R_2^{\mathrm{sat}}/\mathcal R_2$ の像で生成される ideal である。
計算した Smith 因子により、この次数二の商は $k^3$ で $z$ により消える。
$z$ は中心的なので、それが生成する全 ideal も $z$ で消える。
\end{proof}

\begin{proposition}[飽和二次族の中心商と原始的 potential]
\label{prop:nfm-boundary-saturated-active-potential}
命題\ref{prop:nfm-boundary-saturated-quadratic-flat}の平坦二次代数を $B_R$ と書く。
次数一の中心を $U_R=Z(B_R)_1$ とする。
このとき $U_R\subset(B_R)_1$ は自由階数 $r=n-1$ の直和因子であり、
その形成は任意の $R$-底変更に両立する。
その基底は中心正則列をなし、
\[
 \overline B_R=B_R/(U_R),\qquad V_R=(B_R)_1/U_R
\]
は各次数が有限自由な三生成元二次代数である。
一般 fibre は元の active 代数、閉 fibre は上の三角形 Sklyanin 代数となる。

さらに $\overline{\mathcal R}_2=\ker(V_R^{\otimes2}\to(\overline B_R)_2)$ と置くと、
\[
 \mathcal P_R=
 (\overline{\mathcal R}_2\otimes V_R)
        \cap(V_R\otimes\overline{\mathcal R}_2)
       \subset V_R^{\otimes3}
\]
は階数一の直和因子で、底変更に両立する。
元の active 代数の cyclic potential を一般 fibre の入力とすると、
$\mathcal P_R$ は cyclic invariant tensor からなる。
自然な交代商写像
$\operatorname{Alt}:V_R^{\otimes3}\to\bigwedge^3V_R$ は
$\mathcal P_R\xrightarrow{\sim}\bigwedge^3V_R$ を与える。
primitive な potential をこの写像で規格化すると、その symmetric 部分が定める
平面三次は $R$ 上で滑らかである。
これは点曲線 $V(J_R)$ ではなく、active 代数の potential に付随する三次である。
\end{proposition}

\begin{proof}
中心については有限自由加群間の写像
\[
 c:(B_R)_1\longrightarrow
     \operatorname{Hom}_R((B_R)_1,(B_R)_2),\qquad c(a)(b)=ab-ba
\]
を考える。一般 fibre の kernel は元の中心一次空間で次元 $r$ である。
閉 fibre の active 三角形代数には中心一次元がない。
実際、二次部分の基底を
$X^2,Y^2,Z^2,XY,YZ,ZX$ と取ると、
$[\alpha X+\beta Y+\gamma Z,X]$ の $XY,ZX$ 係数は
それぞれ $-\beta,\gamma$ であり、これが零なら $\beta=\gamma=0$ である。
さらに $Y$ と可換なら $\alpha=0$ となる。
従って閉 fibre の中心一次空間も、残りの多項式変数からなる次元 $r$ の空間である。
$c$ の一般・閉 fibre の rank は共に三であり、三次小行列式の一つは単元である。
一般 fibre で rank が三なので四次小行列式は $R$ 上でも全て零である。
単元小行列による消去から、kernel は直和因子で任意底変更に両立する。

$U_R$ の基底を選び、閉 fibre の中心変数の基底へ対応させる。
中心元 $u$ の閉 fibre が正則なら、$ua=0$ から $a\bmod z=0$ となる。
各次数の $R$-自由性を用いてこれを繰り返すと、
$a\in\bigcap_{j\ge0}z^j(B_R)_d=0$ である。
また $zb=ua$ なら、閉 fibre の正則性から $a=za_1$ と書け、
$z$ が正則なので $b=ua_1$ となる。
従って $B_R/(u)$ にも $R$-torsion はない。
各次数は有限なので、この商も次数ごとに有限自由である。
特殊代数では中心変数が正則列をなすため、この議論を順に適用できる。
これで中心正則列と商の全次数の自由性を得る。
構成した短完全列の商は $R$-flat なので、この正則列の計算も底変更と両立する。

$V_R$ の階数は三であり、二次商の階数は六である。
従って $\overline{\mathcal R}_2$ は自由階数三の直和因子である。
三次商は自由階数十なので、
$I_3=\ker(V_R^{\otimes3}\to(\overline B_R)_3)$ は自由階数十七の直和因子である。
二次 presentation から
\[
 (\overline{\mathcal R}_2\otimes V_R)\oplus
 (V_R\otimes\overline{\mathcal R}_2)\longrightarrow I_3,
 \qquad(a,b)\longmapsto a-b
\]
は全射で、その kernel は交わり $\mathcal P_R$ に同型である。
始域の階数は十八なので kernel は自由階数一で、全ての底変更に両立する。
その $V_R^{\otimes3}$ への包含も各 fibre で単射なので primitive である。
一般 fibre の cyclicity は、この tensor 空間の $R$-torsion-free 性により
$R$ 上へ延びる。

閉 fibre の関係を
$r_X=ZY+(p/2)X^2$、$r_Y=XZ+(p/2)Y^2$、
$r_Z=YX+(p/2)Z^2$ とする。
非零 tensor
\[
 w_0=Xr_X+Yr_Y+Zr_Z
     =r_XX+r_YY+r_ZZ
     =XZY+YXZ+ZYX+\frac p2(X^3+Y^3+Z^3)
\]
は交わりを生成し、
\[
 \operatorname{Alt}(w_0)=-3X\wedge Y\wedge Z,\qquad
 \operatorname{Sym}(w_0)
 =3XYZ+\frac p2(X^3+Y^3+Z^3)
 =\frac32 g_{p,2}
\]
である。ここで $\operatorname{Sym}$ も自然な symmetric 商写像とする。
$3$ は単元なので交代射は $R$ 上でも同型である。
例えば $\operatorname{Alt}(w_R)=X\wedge Y\wedge Z$ と規格化し、
$g_R=-2\operatorname{Sym}(w_R)$ と定義すれば、特殊三次は厳密に
\[
 g_R\bmod z=g_{p,2}=\frac p3(X^3+Y^3+Z^3)+2XYZ
\]
となる。
この三次が特異なら、勾配の共通非零零点で全座標が非零となり、
三つの勾配等式を掛けると $p^3=-8$ を得る。
本係数体では $p(p^3+8)$ は単元なので不可能である。
従って判別式は $R$ の単元となり、potential 三次は $R$ 上で滑らかである。
\end{proof}

\begin{proposition}[飽和二次関係に対する三次 class の最初の項]
\label{prop:nfm-boundary-first-cubic-class}
命題\ref{prop:nfm-boundary-first-jet-products}の相対模型において、指定 MF 周期を identity とし、
次数 $(2,-3)$ の cochain $\operatorname{diag}(A,B)$ からの射影を $A_{11}$ とする。
pair homotopy に上の正規化を用いた transfer では、
一次生成元上の三次テンソルの最初の項は
\[
 m_3^*(t^*)=\frac{\varepsilon}{3p}\,acb+O(\varepsilon^2).
\]
符号は補題\ref{lem:nfm-tate-cubic-relation-normalization}の
source 二拡大が $+\Omega$ を回復する規約に合わせた。
$\varepsilon$ で割った三次 class は、命題\ref{prop:nfm-boundary-first-jet-products}の飽和した二次関係による商で非零である。
従って、この飽和した模型における三次 relation の最初の消失次数は一である。
\end{proposition}
\begin{proof}
次数 $(2,-3)$ の boundaries は零、cycles は identity のみである。
従って $A_{11}=0$ という補空間を取れば、指定した射影は相対 contraction に延長できる。
積の homotopy の零次項は $(c,b)$ に対する
$T_{cb,0}=\begin{psmallmatrix}0&0\\-b_0E_{31}&0\end{psmallmatrix}$ だけである。
命題\ref{prop:nfm-separated-boundary-active-products}の
$\phi_iT=T\phi_i=0$ により、零次の三次 tree は全て零である。
従って $\Theta_0=0$ は商の class だけでなく tensor 自体として成立する。
また全ての $C_{i,0}$ の第一行は零である。
従って一次の $A_{11}$ 成分に寄与するのは
\[
 b_0\,(C_{i,1}(-E_{31}))_{11}
 \qquad\text{for }(j,k)=(c,b)
\]
だけである。命題\ref{prop:nfm-boundary-first-jet-products}の行列から、$i=a$ の場合だけ
$-b_0/(pd)=-1/(3p)$ となる。
この規約で寄与する tree の上 block は $-\phi_iT_{jk}$ なので、
主張の正符号を得る。規約の校正は一変数の
\[
 \mathsf D_x=\begin{pmatrix}0&x\\x^2&0\end{pmatrix},\qquad
 \phi=\begin{pmatrix}0&1\\-x&0\end{pmatrix},\qquad
 h=\begin{pmatrix}0&0\\-1&0\end{pmatrix}
\]
で行える。$\phi^2=\delta h$ であり、閉じた三次 tree
$-(h\phi+\phi h)=I_2$ は $+x^3$ を回復する。
この議論は同じ contraction の二つの有限 tree を計算しており、
独立に選んだ二次・三次積を混用していない。

三次 class の非零性は次の双対証明書でも確認できる。
三次語への線型汎関数を
\[
 \ell(aaa)=-p,\quad
 \ell(acb)=\ell(bac)=\ell(cba)=1,\quad
 \ell(bbb)=-p^{-2},\quad\ell(ccc)=-8p^{-2}
\]
とし、その他の語では零とする。
これは三つの二次関係の全ての左右一次倍を消し、$\ell(acb)=1$ である。
従って $[acb]\ne0$ である。
gauge による一次の変化は、一次二次関係の左右一次倍に入るので、
この商での非零性は gauge に依存しない。
\end{proof}

\begin{lemma}[滑らかな potential 三次による active 表示同型の延長]
\label{lem:nfm-boundary-potential-isom-extension}
$R$ を剰余体が標数零の DVR、$K$ をその商体とする。
$V,V'$ を階数三の自由 $R$-加群、
$Q\subset V^{\otimes2}$、$Q'\subset(V')^{\otimes2}$ を階数三の部分束とし、
$B=T_R(V)/(Q)$、$B'=T_R(V')/(Q')$ と置く。
$B_3,B'_3$ は自由階数十とする。
両一般 fibre の cubic syzygy line は cyclic tensor からなり、
両特殊 fibre は $S(0,-2,-p)$ と同型であると仮定する。
ここで剰余体において $p(p^3+8)\ne0$ とする。

このとき、任意の実際の次数付き $K$-同型
$f_K:B'_K\xrightarrow{\sim}B_K$ の次数一部分の射影類は
$R$ 上へ一意に延びる。その線型持上げ $f_R:V'\xrightarrow{\sim}V$ を選ぶと、
$f_R$ は両二次表示の $R$-同型であり、ある $\lambda\in K^\times$ に対して
$f_K|_{V'_K}=\lambda f_R|_{V'_K}$ となる。
\end{lemma}
\begin{proof}
まず
\[
 W=(Q\otimes V)\cap(V\otimes Q)\subset V^{\otimes3}
\]
を取る。
\[
 0\to W\to(Q\otimes V)\oplus(V\otimes Q)
   \to V^{\otimes3}\to B_3\to0
\]
と $B_3$ の自由性から、$W$ は自由階数一で、特殊 fibre と一般 fibre への
底変更に両立する。$Q$ が部分束なので二つの tensor 部分加群は飽和しており、
その交わり $W$ も飽和している。$W'$ も同様である。

$W_K$ 上の三巡回作用が恒等であるという仮定は、$V^{\otimes3}$ の
$R$-torsion-free 性により $W$ 上でも成立する。標数零での自然な分解
\[
 (V^{\otimes3})^{C_3}=\operatorname{Sym}^3V\oplus\bigwedge^3V
\]
を用いる。特殊 fibre $S(0,-2,-p)$ の cyclic potential は
\[
 w_0=-2(yxz+zyx+xzy)-p(x^3+y^3+z^3)
\]
である。その alternating 成分は
\[
 \varepsilon=xyz+yzx+zxy-xzy-zyx-yxz
\]
であって非零である。従って $W\to\bigwedge^3V$ は同型であり、
volume を固定すれば、alternating 成分が $\varepsilon$ となる生成元 $w$ を取れる。
同様に $w'$ を取る。

可換化写像 $\operatorname{ab}:V^{\otimes3}\to\operatorname{Sym}^3V$ により
$g=-\operatorname{ab}(w)/3$ と定めると、その特殊 fibre は適当な線型座標で
\[
 g_0=\frac p3(x^3+y^3+z^3)+2xyz
\]
の非零 scalar 倍となる。これは $p(p^3+8)\ne0$ により滑らかである。
従って $V(g)$ は $R$ 上の滑らかな平面三次であり、$g'$ も同様である。

$f_K$ は内在的に定まる $W'_K$ を $W_K$ へ移す。
alternating 成分を比較すると
\[
 f_K(w')=\det(f_K)w,\qquad f_K(g')=\det(f_K)g
\]
である。従って $f_K$ の射影類は、この二つの滑らかな偏極平面三次の同型である。
Benoist の Theorems~1.6(ii), 1.7 および Lemma~2.2(iii)
\cite[pp.~140--141]{NFMBenoistSeparation} により、
標数零の滑らかな平面三次の射影同型空間は有限であり、指定された一般点の
射影同型は $R$ 上へ一意に延びる。
ここで Theorem~1.7 の separatedness は対角の proper 性であり、
Theorem~1.6(ii) は $3$ を可逆にした基底上での Deligne--Mumford 性を与える。
三角形である点曲線の同型空間へ、この定理を適用しているのではない。

$\operatorname{Pic}(R)=0$ なので、得られた $\operatorname{PGL}_3(R)$ の点を
$f_R\in\operatorname{GL}_3(R)$ へ持ち上げられる。
$Q=Q_K\cap V^{\otimes2}$ と $Q'=Q'_K\cap(V')^{\otimes2}$ により、
一般 fibre での関係保存は $f_R(Q')=Q$ を与える。
従って $f_R$ は次数付き代数の同型である。
\end{proof}

\begin{lemma}[active 中心商での周期付値の評価]
\label{lem:nfm-boundary-active-lattice-valuation}
前補題の同型を用いる。$\Omega\in B_3$ は primitive とし、
$\Omega'\in B'_3$ は
\[
 f_K(\Omega')=a_K\Omega
\]
を満たすとする。$\Omega'$ の $B'_3$ における付値を $e$ とすると
\[
 [a_K]=[z^e u_R],\qquad u_R\in R^\times,
 \quad v_R([a_K])=e\pmod3.
\]
これは full 表示の同型の延長を要せず、実際の周期評価式を active 中心商へ
移せる場合に適用できる。
\end{lemma}
\begin{proof}
$f_K=\lambda f_R$ の次数三部分は $\lambda^3 f_R$ である。
$f_R$ は自由格子の同型なので、$f_R(\Omega')$ の付値も $e$ である。
$\Omega$ が primitive だから $f_R(\Omega')=z^e u_R\Omega$、$u_R\in R^\times$。
従って $a_K=\lambda^3z^e u_R$ である。
\end{proof}

\begin{lemma}[三角形 Sklyanin の自己同型と残余 scalar]
\label{lem:nfm-triangle-sklyanin-central-scalar}
$k$ は $\mu_3$ を含む標数零の体、$s=-p/2$、$p(p^3+8)\ne0$ とし、
\[
 S_0=k\langle X,Y,Z\rangle/(ZY-sX^2,\ XZ-sY^2,\ YX-sZ^2)
\]
と置く。この代数の射影次数付き自己同型群は $\mu_3^2$ であり、
任意の次数付き自己同型は $G=X^3=Y^3=Z^3$ を $k^{\times3}$ 倍する。
\end{lemma}
\begin{proof}
関係行列の determinant は $(1-s^3)XYZ$ であり、三角形上で階数は二である。
従って点対応は三成分を巡回する自己同型の graph である。
次数付き自己同型は三角形を保つので、その線型表示は monomial である。
成分置換は三巡回と可換でなければならず、従って巡回置換に限られる。
対角表示 $\operatorname{diag}(\alpha,\beta,\gamma)$ は
\[
 \beta\gamma=\alpha^2,\qquad
 \gamma\alpha=\beta^2,\qquad
 \alpha\beta=\gamma^2
\]
を満たす。射影的に $\alpha=1$ とすると $\beta^3=1$、$\gamma=\beta^2$ である。
従って射影群は、$\operatorname{diag}(1,\zeta,\zeta^2)$ と巡回置換による
位数九の群である。これらの線型持上げは $G$ を固定する。
残る grading scalar $\lambda$ は $G$ を $\lambda^3G$ へ移す。
\end{proof}

\begin{proposition}[全次元の完備局所周期位相]
\label{prop:nfm-boundary-period-local-kummer}
元の一ブロック一般族、正準半古典 frame、および固定 MF・window・
一般点生成性・coherent 飽和性の入力の下で、$n\ge3$ とする。
$k_0=\mathbb C(p,\mathbf m)$、$R=k_0[[z]]$、$K=k_0((z))$、$q=2+z$ と置く。
ここで $\mathbf m$ は任意個数の独立な混合・純中心係数である。
$D_R=D(r(p,2+z))$ とすると、正準強化点型 class の周期位相は
\begin{equation}\label{eq:nfm-boundary-period-local-kummer}
 \left.\gamma_{\mathrm{sc},n}\right|_K=[D_R]
   \quad\text{in }K^\times/K^{\times3},\qquad
 D_R=-\frac{(p^3+8)^2}{27}z+O(z^2).
\end{equation}
従ってこの torsor の局所体は $K(\sqrt[3]{D_R})$ であり、
$D=0$ の一般点での mod-three 付値は一、tame 分岐指数は三である。
この結論は全ての $n\ge3$ で同じである。
\end{proposition}

\begin{proof}
\emph{復元した三次の active 格子への移行。}
中心座標の pole 消去を $F_R\equiv J_R\pmod{z^2}$ となるまで行う。
これは一般 fibre 上の指定方程式と実際の周期を保つ座標変更であり、
体拡大や方程式生成元の自由な再規格化ではない。
命題\ref{prop:nfm-period-phase-low-products}により、一般 fibre の実際の
次数付き表示同型 $f_K$ に対して
\[
 \gamma_{\mathrm{sc},n}|_K=[a_K],\qquad
 f_K(\Omega_p)=a_K\Omega
\]
である。この式を次数一中心による商へ移す。
表示同型は内在的な中心を保ち、source 中心三次の active 像は非零なので、
同じ scalar $a_K$ を得る。

命題\ref{prop:nfm-boundary-saturated-quadratic-flat}と
\ref{prop:nfm-boundary-saturated-active-potential}により、target の飽和した
active 二次表示 $\overline B'_R$ は全次数で自由である。
ここで三次テンソルを $\Theta_R=m_3^*(t_p^*)$ と書く。
分離した零次模型では、active 三次 tree は文字通り零であり、
中心入力を含むものも厳密単位性により零である。
従って \emph{二次関係による商を取る前から}
$\Theta_R=z\omega_R$ と書ける。
中心商への射影を $\pi_R$ とすると
\[
 \left(\frac{\pi_R[\Theta_R]}z\right)\bmod z
      =\pi_0[\omega_R\bmod z].
\]
中心部分束の一次の傾きや、飽和関係の次の jet はこの式に寄与しない。
補題\ref{lem:nfm-boundary-full-active-first-jet}と
命題\ref{prop:nfm-boundary-first-cubic-class}から
\[
 \overline\Omega_p=z\omega_R,\qquad
 \omega_R\bmod z=\frac1{3p}[acb]\ne0.
\]
従って飽和 active 格子における $\overline\Omega_p$ の付値は正確に一である。

\emph{source 中心三次の原始性。}
source active 代数も同じ $R$ 上の Hesse 表示を持つ。
一般 fibre と特殊 fibre $S(0,-2,-p)$ の Hilbert 級数が共に
$(1-t)^{-3}$ なので、上と同じ次数別の議論で全次数が自由である。
その cyclic potential の可換部分は
$g_{p,q}=p(X^3+Y^3+Z^3)/3+qXYZ$ である。
量子化論文\cite[Hesse 型の普遍混合中心三次式の定理]{CubicQuantization}の
厳密な中心公式を $\hbar'=p$、$t=q/p$ で評価し、全体を $p$ 倍すると、
中心変数を零とした像は
\[
 \overline\Omega=
 \frac{pU}{3}(X^3+Y^3+Z^3)+V\operatorname{sym}(XYZ),
\]
\[
 U=1+\frac{q^2}{2}+\frac{p^3q}{6}-\frac{q^4}{48},\qquad
 V=\frac{(4-q^2)(12q+q^3-8p^3)}{48}.
\]
$q=2$ では $U=(p^3+8)/3$、$V=0$ であり、
$G=X^3=Y^3=Z^3\ne0$ を用いて
\begin{equation}\label{eq:nfm-boundary-source-center-primitive}
 \overline\Omega\bmod z=\frac{p(p^3+8)}3G
\end{equation}
となる。従って source 三次は primitive である。
ここでは固定入力の中心公式を用いており、その中心性全体の独立な再証明を
今回の一次 jet の検算で置き換えてはいない。

\emph{指定一般点同型の延長と単元の比較。}
両 active 表示の特殊 fibre は $S(0,-2,-p)$ であり、その potential 三次は滑らかである。
補題\ref{lem:nfm-boundary-potential-isom-extension}を実際の $f_K$ に適用し、
$f_K=\lambda f_R$ となる格子同型 $f_R$ を得る。
次数三では $\lambda^3$ が作用するので、Kummer 類には影響しない。
full 表示の中心変数方向の同型が $R$ 上へ延びることは必要ない。

特殊 fibre の基準同型を
$a\mapsto X$、$b\mapsto Y/p$、$c\mapsto2Z/p$ とする。
$XZY=-(p/2)X^3$ より $acb\mapsto-G/p$ であり、
\[
 f_0(\omega_R\bmod z)=-\frac1{3p^2}G.
\]
$f_R\bmod z$ とこの基準同型との差は source 特殊代数の次数付き自己同型である。
補題\ref{lem:nfm-triangle-sklyanin-central-scalar}により、
その $G$ 上の scalar は $k_0^{\times3}$ に属する。
従って補題\ref{lem:nfm-boundary-active-lattice-valuation}と
式\eqref{eq:nfm-boundary-source-center-primitive}から
\[
 [a_K]=[zu_R],\qquad
 u_R\bmod z=-\frac1{p^3(p^3+8)}
                  \quad\bmod k_0^{\times3}.
\]
一方、$D_R/z$ の剰余との比は
\[
 \frac{-1/[p^3(p^3+8)]}{-(p^3+8)^2/27}
       =\left(\frac3{p(p^3+8)}\right)^3.
\]
$R$ は完備、$3$ は可逆なので、$1+zR$ の任意の元は $1+zR$ 内で一意な三乗根を持つ。
従って $[a_K]=[D_R]$ である。
付値一から Kummer 拡大は次数三で全分岐し、剰余標数零なので tame である。
さらに補題\ref{lem:nfm-period-hesse-eight-boundary-lines}で
Hesse 写像の $q=2$ における分岐指数は一なので、元の $D$-付値も一 mod~三となる。
\end{proof}

\begin{corollary}[正準位相体の次数と大域式への残る条件]
\label{cor:nfm-boundary-phase-degree-three}
標準 split 一般係数体 $k_n$ 上で、全 $n\ge3$ の正準強化点型の周期位相は非自明であり、
その $\mu_3$-torsor は連結な三次体 $M/k_n$ である。
既存の scalar gerbe を $\alpha$ とすると、指定周期を標識した実現の最小次数は
\[
 3\operatorname{ind}(\alpha_M)\in\{3,9\}.
\]
$n=3,4$ では、この式は全 full class に共通である。

さらに $\overline\gamma_n$ の torsor が
$\operatorname{Spec}\mathbb C[s,t,1/(\Delta D)]$ 上で有限 étale に延びるという
追加条件の下では、大域的に $\overline\gamma_n=[D]$ となる。
この追加条件と $\operatorname{ind}(\alpha_M)$ の値は本計算では未決定である。
\end{corollary}
\begin{proof}
体拡大後に非自明なので、元の位相も非自明である。
$\mu_3\subset k_n$ かつ三が素数なので、対応する有限 étale 三次代数は体である。
標識付き最小次数の既存公式 $[M:k_n]\operatorname{ind}(\alpha_M)$ を適用する。
元の scalar index は三なので、拡大後の index は一または三である。
低次元での class 不変性は既存の周期整合定理による。
追加の非分岐条件の下では補題\ref{lem:nfm-period-two-boundary-criterion}により
候補が $1,[D],[D]^2$ に限られ、今回の $D$-付値一が $[D]$ を選ぶ。
\end{proof}

\begin{remark}[局所式の範囲と異なる二つの被覆]
正確な Kummer 等式\eqref{eq:nfm-boundary-period-local-kummer}は完備 trait 上の主張である。
未完備 trait 上でも mod-three 付値は一だが、剰余が三乗である単元が
その商体で三乗となることまでは主張しない。
また命題\ref{prop:nfm-hesse-boundary-third-division}の三分点 torsor は、
この局所体の三つのコピーからなる九次被覆である。
正準位相の三次被覆とは、一成分の体が一致するだけであり、torsor 全体は異なる。
全次元での局所式の一致から、二変数不変量体上での位相の一致や
他の divisor での非分岐性を結論してはいない。
任意の方程式 frame の変更則や独立標的尺度による次数九の例も引き続き有効である。
\end{remark}
